
\section{РЕЗУЛЬТАТЫ ЭКСПЕРИМЕНТОВ И АНАЛИЗ ДАННЫХ}

\subsection{Условия проведения эксперимента}

% TODO: подставить реальное число сгенерированных заявок и условия эксперимента
В рамках экспериментальной проверки работы имитатора было сгенерировано
и поставлено в очередь диспетчера $N = $ \texttt{<N>} заявок. Все
четыре вычислительных узла кластера на момент начала эксперимента
находились в состоянии готовности (\texttt{idle}).

\subsection{Результаты постановки заявок в очередь}

% TODO: вставить скриншот/фрагмент вывода squeue в процессе выполнения
% \begin{figure}[H]
% \centering
% \includegraphics[width=0.9\textwidth]{images/squeue_output.png}
% \caption{Состояние очереди заявок диспетчера SLURM в процессе эксперимента}
% \end{figure}

По завершении выполнения всех сгенерированных заявок были собраны
итоговые статистические данные, фрагмент которых приведён в
таблице~\ref{tab:results}.

% TODO: подставить реальные данные из sacct (5-10 характерных строк)
\begin{table}[H]
\caption{Сравнение плановых и фактических характеристик выполнения заявок (фрагмент)}
\label{tab:results}
\centering
\begin{tabular}{|c|c|c|c|c|}
\hline
JobID & $T_{\text{заявки}}$, с & $T_{\text{факт}}$, с & NNodes & State \\
\hline
\texttt{<...>} & \texttt{<...>} & \texttt{<...>} & \texttt{<...>} & COMPLETED \\
\hline
\end{tabular}
\end{table}

\subsection{Анализ отклонений «план -- факт»}

Сопоставление запрошенного пользователем времени выполнения
$T_{\text{заявки}}$ и фактического времени выполнения $T_{\text{факт}}$
по результатам эксперимента позволило вычислить выборочные значения
среднего и стандартного отклонения логарифма отношения $\tau$ для
синтетических данных и сравнить их с параметрами $\mu$ и $\sigma$,
полученными на этапе анализа исходного набора \texttt{acct\_0921-0923}
(подраздел~2.4).

% TODO: подставить реальные итоговые значения и их сравнение с исходными mu/sigma
Полученные по результатам эксперимента значения
$\hat{\mu} = $ \texttt{<значение>} и $\hat{\sigma} = $ \texttt{<значение>}
\texttt{<близки / отличаются от>} соответствующих параметров исходной
выборки, что подтверждает (либо позволяет уточнить) адекватность
выбранной статистической модели генерации пользовательских заявок.

\subsection{Выводы по результатам экспериментов}

Проведённый эксперимент показал, что разработанный имитатор
пользовательских заявок обеспечивает корректную автоматическую
постановку заданий в очередь диспетчера SLURM, а статистические
характеристики сгенерированной нагрузки воспроизводят характер
распределения времени выполнения и ошибки его оценки, наблюдаемый в
исходном наборе учётных данных. Это подтверждает применимость
предложенного подхода к моделированию пользовательской нагрузки на
вычислительный кластер при отсутствии доступа к данным реальной
эксплуатационной системы в реальном масштабе времени.